\chapter{Descrizione dello stage}
\label{chap:descrizione-stage}

\section{Introduzione al progetto}
\noindent Il progetto consiste nello sviluppo di un'applicazione multipiattaforma per le note spesa, progettata per supportare l'ufficio amministrativo nelle attività quotidine. L'applicazione mobile dovrà essere cross-platform, obiettivo che sarà raggiunto utilizzando il framework \gls{dotnetmauig} di Microsoft, che consente di scrivere l'interfaccia e la logica di business usando \gls{csharpg} e \gls{xamlg}, permettendo di creare applicazioni per iOS e Android con un unico codice base. La comunicazione tra l'applicazione mobile e il sistema centrale avverà tramite \gls{apirestg}, garantendo così un'integrazione fluida ed efficiente con l'infrastruttura esistente. 


\section{Contesto aziendale}
\noindent Durante il periodo di stage ho avuto modo di interagire costantemente con i membri del team di lavoro, in particolare con il tutor aziendale, il quale si è sempre dimostrato disponibile a fornire supporto e chiarimenti riguardo alle attività da svolgere e ai dubbi relativi al progetto. 
Il confronto con il tutor avveniva principalmente in presenza, attraverso incontri periodici in cui venivano monitorati l'andamento del lavoro e le eventuali difficoltà incontrate. In tali occasioni venivano inoltre pianificate le attività successive, definendo in modo chiaro gli obiettivi da raggiungere. 
Oltre agli incontri programmati, era possibile confrontarsi con il tutor per eventuali chiarimenti, garantendo così un supporto costante durante lo svolgimento delle attività.
Il supporto del tutor si è rivelato utile per comprendere le dinamiche aziendali e le aspettative del progetto, favorendo un inserimento più consapevole nel contesto lavorativo. 


\section{Analisi preventiva dei rischi}
\noindent Durante la fase di analisi iniziale sono stati individuati alcuni rischi potenziali legati allo sviluppo del progetto e all'adozione di nuove tecnologie.
Sono state quindi elaborate delle possibili soluzioni per far fronte a tali criticità.
\begin{risk}{Inesperienza con le tecnologie adottate}
    \riskdescription{il progetto prevede l'utilizzo di \gls{csharpg} e \gls{xamlg}, \gls{dotnetmauig}, oltre all'editor Visual Studio. Tali tecnologie non erano state precedentemente utilizzate e ciò potrebbe rallentare i tempi di implementazione}
    \risksolution{si è deciso di dedicare una fase iniziale all'apprendimento tramite documentazione ufficiale e tutorial, affiancata alla realizzazione di piccoli progetti di prova al fine di acquisire familiarità e competenze necessarie prima dello sviluppo del progetto}
    \label{risk:tecnologie-nuove} 
\end{risk}
\begin{risk}{Gestione delle scadenze e pianificazione}
    \riskdescription{il progetto potrebbe subire ritardi a causa di una stima non corretta dei tempi o delle risorse necessarie, oppure per l'insorgere di imprevisti tecnici o bug complessi da risolvere}
    \risksolution{le attività sono state pianificate in modo realistico e incrementale, monitorando costantemente l'avanzamento e prevedendo margini temporali per la gestione degli imprevisti. Sono state inoltre effettuate revisioni periodiche per individuare e risolvere tempestivamente eventuali criticità}
    \label{risk:scadenze-pianificazione} 
\end{risk}


\section{Pianificazione}
\begin{table}[H]
    \centering
    \rowcolors{2}{}{tableGray}
    \begin{tabularx}{\textwidth}{|M{2cm}|M{1cm}|Y|}
        \hline
        \textbf{Settimana} & \textbf{Ore} & \textbf{Attività}\\ 
        \hline
        1 & 32 & Introduzione al progetto e alle tecnologie da utilizzare. Analisi iniziale dei requisiti e delle funzionalità richieste del sistema. \\
        \hline
        2 & 40 & Studio delle tecnologie \gls{csharpg}, \gls{xamlg} e del framework \gls{dotnetmauig}, con approfondimento del pattern \gls{mvvmg}. Studio delle metodologie di sviluppo delle API REST e prima stesura dell'analisi dei requisiti funzionali e non funzionali. \\
        \hline
        3 & 40 & Conclusione della fase di studio e completamento dell'analisi dei requisiti. Progettazione dell'architettura dell'applicazione e avvio dello sviluppo. \\
        \hline
        4 & 40 & Sviluppo delle principali funzionalità dell'applicazione mobile e implementazione dell'interfaccia grafica. \\
        \hline
        5 & 40 & Prosecuzione dello sviluppo dell'applicazione e integrazione delle funzionalità implementate. \\
        \hline
        6 & 32 & Sviluppo dei web services per la gestione del flusso di dati tra l'ERP aziendale e l'applicazione mobile. \\
        \hline
        7 & 40 & Implementazione delle procedure di integrazione con il gestionale, attività di testing e validazione del progetto. \\
        \hline
        8 & 40 & Redazione della documentazione tecnica del sistema e del manuale utente. \\
        \hline
    \end{tabularx}
    \vspace{8pt}
    \caption{Pianificazione del lavoro}
    \label{tab:pianificazione}
\end{table}

\newpage